\scanpage{112}
\mspage{27}
% archive sheet number: 111
\deleted{Notons que}

\[
\deleted{\inserted{H^i}\quad H^{i-2}(D'_r,\Lambda)\otimes\mu_n(k)^{\otimes-1}\longrightarrow H^i}
\]

\deleted{$\unsure{H^i|\Lambda}$} Je dis que l'hom. $\alpha_i$ de (4.5.5) est \uline{nul}. En effet, les calculs explicites précédents (liés à l'hyp. de récurrence) prouvent que

\[
H^{\unsure{*|\alpha}}(U',\deleted{\unsure{E}}F)\unsure{{-}|\longrightarrow}H^*(D'_r,F)
\]

est \deleted{\unsure{surj.}} \inserted{un isom.} si $F$ est loc. constant et annulé par $n$, p.ex. si $F=\mu_n^{\otimes-1}$. Il suffit donc de prouver que le composé

\[
H^{i-2}(U',\deleted{\unsure{\mu_n^{\otimes-1}}}\Lambda(-1))\longrightarrow H^{i-2}(D'_r,\deleted{\unsure{\mu_n}}\Lambda(-1))\xrightarrow{\alpha_i}H^i(U',\Lambda)
\]

est nul. Or en vertu de 4.3, \deleted{\unsure{c'est le}} composé est le cup-produit par la classe \deleted{de} $\gamma_{U'}(D'_r)\inserted{=c^1(\underline{\unsure{\mathcal O|O}}_{U'}(D'_r))}$, \deleted{\unsure{or|si} $D'_r$ est un diviseur \unsure{p}} cette classe est induite par la classe \deleted{\unsure{de} $\unsure{D_r}$} $\gamma_X(D_r)$, qui est nulle puisque $X$ est str. local (ou \unsure{parce} que $D_r$ est principal). Donc la suite exacte (4.5.5) se décompose en suites exactes courtes

\[
0\longrightarrow H^i(U',\Lambda)\longrightarrow H^i(U,\Lambda)\longrightarrow H^{i-1}(D'_r,\Lambda)\deleted{\unsure{\otimes\mu_n(k)^{-1}}}(-1)\longrightarrow0.
\tag*{(4.5.6)}
\]

% Cancelled annotations below the first and third cohomology terms of (4.5.6).
\[
\deleted{\underset{\left(\bigwedge^i(\Lambda^{r-1})\right)\otimes\mu_n(k)^{\otimes-i}}{\wr}\qquad\underset{\bigwedge^{i-1}}{\wr}}
\]

% The final paragraph is struck through and enclosed by cancellation lines.
\deleted{D'autre part, la décomposition de \deleted{\unsure{\(\Lambda\) de la fibre}} $\Lambda^r=\coprod_{1\leq i\leq\unsure{r|r-1}}\Lambda_{D_i}$ suivant \deleted{$\unsure{\Lambda}$} $\Lambda^{r-1}\times\Lambda$, où le dernier facteur $\Lambda$ est $\Lambda_{D_r}$, donne}
